\documentclass[12pt]{article}
\usepackage[utf8]{inputenc}
\usepackage{amsmath}
\usepackage{amssymb}
\usepackage{geometry}
\usepackage{multicol}
\geometry{a4paper, margin=0.75in}

\title{Mathematics Test - Two Column Layout}
\author{Generated by Image to LaTeX Pipeline}
\date{\today}

\begin{document}

\maketitle

\begin{multicols}{2}

\section*{Questions}

\subsection*{Q1}
Solve: $3x - 7 = 14$

\textbf{Options:}
\begin{itemize}
    \item[A)] $x = 7$
    \item[B)] $x = 5$
    \item[C)] $x = 3$
    \item[D)] $x = 21$
\end{itemize}

\subsection*{Q2}
Find $\frac{d}{dx}[\sin(x)]$

\textbf{Options:}
\begin{itemize}
    \item[A)] $\cos(x)$
    \item[B)] $-\cos(x)$
    \item[C)] $\sin(x)$
    \item[D)] $\tan(x)$
\end{itemize}

\subsection*{Q3}
Evaluate: $\lim_{x \to 0} \frac{\sin(x)}{x}$

\textbf{Options:}
\begin{itemize}
    \item[A)] $0$
    \item[B)] $1$
    \item[C)] $\infty$
    \item[D)] Undefined
\end{itemize}

\subsection*{Q4}
What is $5!$ (5 factorial)?

\textbf{Options:}
\begin{itemize}
    \item[A)] $25$
    \item[B)] $50$
    \item[C)] $120$
    \item[D)] $720$
\end{itemize}

\subsection*{Q5}
Solve: $\sqrt{x + 3} = 5$

\textbf{Options:}
\begin{itemize}
    \item[A)] $x = 2$
    \item[B)] $x = 22$
    \item[C)] $x = 25$
    \item[D)] $x = 8$
\end{itemize}

\columnbreak

\section*{Solutions}

\subsection*{S1: Answer A}
\textbf{Work:}
\begin{align*}
    3x - 7 &= 14 \\
    3x &= 21 \\
    x &= 7
\end{align*}

\subsection*{S2: Answer A}
\textbf{Explanation:}
The derivative of $\sin(x)$ is $\cos(x)$. This is a fundamental derivative that should be memorized.

\subsection*{S3: Answer B}
\textbf{Explanation:}
This is a famous limit:
$$\lim_{x \to 0} \frac{\sin(x)}{x} = 1$$

This can be proven using L'Hôpital's rule or geometric arguments.

\subsection*{S4: Answer C}
\textbf{Work:}
$$5! = 5 \times 4 \times 3 \times 2 \times 1 = 120$$

\subsection*{S5: Answer B}
\textbf{Work:}
\begin{align*}
    \sqrt{x + 3} &= 5 \\
    x + 3 &= 25 \\
    x &= 22
\end{align*}

\textbf{Check:} $\sqrt{22 + 3} = \sqrt{25} = 5$ ✓

\end{multicols}

\vspace{1em}

\section*{Summary}

\begin{center}
\begin{tabular}{|c|c|c|}
\hline
\textbf{Question} & \textbf{Answer} & \textbf{Concept} \\
\hline
Q1 & A ($x = 7$) & Linear Equations \\
Q2 & A ($\cos(x)$) & Derivatives \\
Q3 & B ($1$) & Limits \\
Q4 & C ($120$) & Factorials \\
Q5 & B ($x = 22$) & Radicals \\
\hline
\end{tabular}
\end{center}

\vspace{1em}

\textbf{Score Calculation:}
\begin{itemize}
    \item Each question: 20 points
    \item Total: 100 points
    \item Passing grade: 60\% (3 correct answers)
\end{itemize}

\end{document}
